% Adaptación del preámbulo de Paper III v1.5 español; no plantilla standalone de Pandoc.
\documentclass[spanish,11pt,a4paper]{article}
\usepackage[margin=1in]{geometry}
\usepackage{fontspec}
\setmainfont{lmroman10-regular.otf}[BoldFont=lmroman10-bold.otf,ItalicFont=lmroman10-italic.otf,BoldItalicFont=lmroman10-bolditalic.otf]
\setsansfont{lmsans10-regular.otf}[BoldFont=lmsans10-bold.otf]
\setmonofont{lmmono10-regular.otf}[Scale=MatchLowercase]
\newlength{\LeanCharWidth}
\usepackage{amsmath,amssymb}
\usepackage[spanish,es-noshorthands,es-tabla]{babel}
\usepackage{microtype,parskip}
\usepackage{graphicx,float,caption}
\usepackage{longtable,booktabs,array,calc}
\usepackage{xcolor}
\usepackage{xurl}
\usepackage{hyperref,bookmark}
\usepackage{etoolbox,needspace}
\setcounter{secnumdepth}{-\maxdimen}
\setlength{\emergencystretch}{3em}
\providecommand{\tightlist}{\setlength{\itemsep}{0pt}\setlength{\parskip}{0pt}}
\newcounter{none}
\AtBeginEnvironment{longtable}{\Needspace{8\baselineskip}\small}
\pretocmd{\section}{\Needspace{5\baselineskip}}{}{}
\pretocmd{\subsection}{\Needspace{4\baselineskip}}{}{}
\pretocmd{\subsubsection}{\Needspace{4\baselineskip}}{}{}
\captionsetup{font=small,labelfont=bf}
\urlstyle{same}
\hypersetup{colorlinks=true,linkcolor=blue,urlcolor=blue,citecolor=blue,pdflang={es-ES}}
\title{De los certificados de realizabilidad a la cota de Erdős 81}
\author{Juan Pablo Traverso Gianini\\
Investigador independiente, Santiago, Chile\\
\href{mailto:jtraverso@gmail.com}{\nolinkurl{jtraverso@gmail.com}}\\
ORCID: 0009-0003-6068-4096}
\date{}\begin{document}\raggedbottom\maketitle
\textbf{Nota de investigación I de la serie, complementaria de los
Papers IV y V --- \mbox{Preprint 1.0}}\\
\textbf{Fecha:} 18 de octubre de 2026\\
Esta nota reorganiza parte de la demostración del Paper IV mediante la
interfaz de transferencia del Paper V. No afirma una nueva solución del
Problema 81 de Erdős ni una demostración independiente. El companion
formal separado identifica el paquete final de fuentes y sus
dependencias.

\subsection{Resumen}\label{resumen}

Reorganizamos parte de la demostración ya establecida de la cota de
partición en cliques del Problema 81 de Erdős como una aplicación del
teorema de transferencia por clases realizables del Paper V, utilizando
el redondeo mixto y las construcciones físicas casi extremales del Paper
IV. El certificado pertinente compara la ganancia de un empaquetamiento
realizado con una demanda de referencia; su pérdida debe ajustarse al
presupuesto de partición del grafo original. Lejos del régimen extremal,
una holgura cuadrática paga el redondeo. Cerca de ese régimen, la
localización por descenso y las cuentas terminales conservadas del Paper
IV proporcionan un empaquetamiento cuya completación ya satisface la
cota entera ajustada. Así se recupera la conclusión de IV, para todos
los grafos cordales suficientemente grandes de orden \(n\), una
partición en cliques de orden a lo sumo cuatro con a lo sumo
\(\left\lfloor n(n+1)/6\right\rfloor\) piezas, y un error aditivo
absoluto para cualquier orden. El argumento de primera entrada de {[}7,
Sección 8{]} y el desarrollo de Okechukwu {[}8{]} son antecedentes
explícitos; aquí se utiliza el descenso de IV. Distinguimos el papel de
la transferencia de esa localización estructural y explicamos por qué no
es necesario transportar ninguna partición hacia atrás a lo largo del
descenso.

\subsection{1. El presupuesto de
partición}\label{el-presupuesto-de-particiuxf3n}

Una partición en cliques de un grafo simple finito \(G\) particiona
\(E(G)\) en los conjuntos de aristas de subgrafos completos de orden al
menos dos. Denotamos por \(\operatorname{cp}(G)\) su tamaño mínimo y por
\(c_4(G)\) el mínimo correspondiente cuando todas las piezas tienen
orden a lo sumo cuatro. Definamos

\[
M(n)=\left\lfloor\frac{n(n+1)}6\right\rfloor.
\tag{1.1}
\]

La pregunta de Erdős, Ordman y Zalcstein {[}1,4{]} pide
\(\operatorname{cp}(G)\le n^2/6+O(n)\) cuando \(G\) es cordal. El Paper
IV {[}2{]} demuestra la cota eventual más fuerte \(c_4(G)\le M(n)\), con
una corrección aditiva absoluta para los órdenes restantes. Aquí
recuperamos esa conclusión mediante el lenguaje de certificados del
Paper V {[}3{]}. El propósito es identificar la interfaz exacta entre la
estructura, la realización física y el número de piezas de la partición.
El problema ya está resuelto en IV y en los otros dos desarrollos
comparados allí: el preprint firmado Anonymous {[}7{]} y Okechukwu
{[}8{]}. Nuestro propósito es expositivo y organizativo. Conservamos los
insumos matemáticos de IV y su atribución, sin ofrecer otra solución
independiente.

\Needspace{4\baselineskip}Para un clique \(Q\), definimos su ganancia mediante

\[
g(Q)=\binom{|Q|}{2}-1.
\]

Si un empaquetamiento \(P\) está formado por cliques disjuntos en
aristas, completar individualmente todas las aristas no cubiertas
produce una partición de tamaño

\[
e(G)-g(P),\qquad g(P)=\sum_{Q\in P}g(Q).
\tag{1.2}
\]

En efecto, cada clique seleccionado reemplaza tantos bloques
individuales como aristas contiene y aporta un bloque en su lugar. La
disjunción en aristas permite sumar estos ahorros sin contar ninguna
arista dos veces. No se modifica ninguna arista exterior a los cliques
seleccionados.

En el modelo mixto utilizado a continuación, las piezas disponibles son
triángulos y copias de \(K_4\), con ganancias dos y cinco. Por tanto,
una estimación de ganancia proporciona una estimación de partición
efectiva sin requerir otro argumento de cobertura. Todos los
empaquetamientos obtenidos en esta nota están formados por cliques del
grafo original.

\subsection{2. Transferencia con un coste
explícito}\label{transferencia-con-un-coste-expluxedcito}

Utilizamos la siguiente forma de la transferencia por clases realizables
del Paper V, con pesos y pérdidas naturales. Tomemos una colección
finita de ocurrencias auxiliares ponderadas, de peso total \(D\),
eliminemos ocurrencias de peso total \(z\) y particionemos las restantes
en \(q\ge1\) clases. Para cada clase \(i\), construyamos un
empaquetamiento físico \(P_i\) en \(G\) y una pérdida no negativa
\(r_i\) tales que

\[
\text{peso de la clase }i\le g(P_i)+r_i.
\tag{2.1}
\]

Las clases no tienen que ser independientes en el grafo auxiliar de
conflictos. La factibilidad se exige a sus realizaciones físicas. Sea
\(C=z+\sum_i r_i\). Sumando (2.1), alguna clase satisface

\[
g(P_i)\ge\frac{D-C}{q}.
\tag{2.2}
\]

Las clases restantes tienen peso total \(D-z\). Al sumar sus
desigualdades de realización se obtiene
\(\sum_i g(P_i)\ge D-z-\sum_i r_i=D-C\). Al menos una de las \(q\)
ganancias es mayor o igual que su promedio. Elegimos ese único
empaquetamiento; no tomamos la unión de los empaquetamientos, que pueden
compartir aristas.

Este es el teorema de transferencia por clases realizables {[}3{]}. Con
\(L=D/q\) y \(\rho=C/q\), la completación da

\[
c_4(G)\le e(G)-L+\rho.
\tag{2.3}
\]

\Needspace{6\baselineskip}
\textbf{Proposición 2.1.} Un certificado cuyas salidas físicas son
triángulos y copias de \(K_4\) demuestra \(c_4(G)\le T\) siempre que

\[
qe(G)+C\le qT+D.
\tag{2.4}
\]

\textbf{Demostración.} Elijamos la salida de (2.2). Su completación
tiene tamaño a lo sumo \(e(G)-D/q+C/q\), que por (2.4) es a lo sumo
\(T\).

Resulta útil comparar con el óptimo fraccional mixto \(W^*(G)\). Si
\(L\le W^*(G)\), definamos

\[
\Delta=M(n)-e(G)+W^*(G),\qquad
\gamma=W^*(G)-L.
\]

La ecuación (2.3) se convierte en

\[
c_4(G)\le M(n)-\Delta+\gamma+\rho.
\tag{2.5}
\]

Para ver la cancelación, sustituyamos \(L=W^*(G)-\gamma\) en (2.3) y
utilicemos después \(e(G)-W^*(G)=M(n)-\Delta\). Las dos pérdidas quedan
medidas respecto del mismo presupuesto: \(\gamma\) es la diferencia
entre el óptimo fraccional y la demanda de referencia por clase,
mientras que \(\rho\) es el coste de realización por clase.

La condición suficiente es \(\gamma+\rho\le\Delta\). Una cota uniforme
de la brecha de integralidad para todo grafo cordal es más fuerte de lo
que requiere esta aplicación. El certificado solo debe ajustarse al
margen del grafo considerado. Escribir (2.5) no proporciona, por sí
solo, ni un defecto de coloreabilidad pequeño ni una realización de bajo
coste.

Para el ensamblaje final utilizamos una clase. Fijemos una referencia
entera \(R\) antes de construir la salida. Si un empaquetamiento físico
\(P\) satisface \(g(P)\le R\), una única demanda auxiliar de peso \(R\)
se realiza mediante \(P\) con pérdida \(R-g(P)\). Su certificado
satisface (2.4) exactamente cuando

\[
e(G)+R-g(P)\le T+R.
\tag{2.6}
\]

La referencia \(R\) se cancela de ambos lados de (2.6). Esta condición
es, por tanto, precisamente \(e(G)-g(P)\le T\), el coste de completación
ya identificado en (1.2). La hipótesis \(g(P)\le R\) garantiza que la
pérdida registrada sea no negativa.

La elección \(R=e(G)\) sirve para cualquier empaquetamiento disjunto en
aristas, pues \(g(P)\le e(G)\). Esta especialización expresa una
construcción física previamente establecida mediante la interfaz de
transferencia; no mejora esa construcción por promediación. El enunciado
general para \(q\) clases sigue siendo útil cuando se dispone de varias
realizaciones, pero la deducción que sigue no necesita una mejora de ese
tipo.

\subsection{3. Los resultados de entrada del Paper
IV}\label{los-resultados-de-entrada-del-paper-iv}

Para precisar las dependencias, enunciamos los dos resultados
combinatorios de entrada independientemente del teorema final de
partición cordal. Sus demostraciones pertenecen al Paper IV. Sea
\(W^*(G)\) la ganancia fraccional mixta máxima, con capacidad uno en
cada arista. Puesto que se trata de un programa lineal racional finito,
el máximo se alcanza con un certificado primal-dual racional.

\Needspace{6\baselineskip}
\textbf{Resultado de entrada 3.1.} Para todo \(\zeta>0\) existe
\(N_\zeta\) tal que todo grafo \(G\) de orden \(n\ge N_\zeta\) tiene un
empaquetamiento físico \(P\) de triángulos y copias de \(K_4\) que
satisface

\[
W^*(G)-g(P)\le\zeta n^2.
\tag{3.1}
\]

La interfaz formal recibe una tolerancia racional positiva. Esto basta
para la tolerancia real enunciada: se elige un racional positivo menor
que \(\zeta\) y se aplica la interfaz a ese número. El error resultante
también es a lo sumo \(\zeta n^2\).

El Resultado de entrada 3.1 es exactamente el caso \(r=4\) del Teorema
3.4 de Rohatgi--Urschel--Wellens {[}11{]}: sus objetivos integral y
fraccional asignan a cada clique la ganancia \(\binom{|Q|}{2}-1\), que
aquí vale dos o cinco. Su pérdida uniforme es \(o(n^2)\). Haxell--Rödl
{[}9{]} y Yuster {[}10, Teorema 1.2{]} son los antecedentes de
transferencia citados allí. El proyecto formal de esta nota utiliza la
implementación demostrada del Paper IV {[}2, Teorema 3.1{]}, construida
con la infraestructura del Paper III; no se afirma un enunciado de
redondeo más fuerte que {[}11{]} para esta aplicación. La pérdida
subcuadrática basta en el régimen lejano y no es una cota de error
lineal. La cordalidad interviene en el segundo resultado de entrada.

\Needspace{6\baselineskip}
\textbf{Resultado de entrada 3.2.} Existe un \(\eta>0\) absoluto tal
que, para todo grafo cordal \(G\) suficientemente grande que satisfaga

\[
e(G)-W^*(G)\ge\frac{n^2}{6}-\eta n^2,
\tag{3.2}
\]

la construcción casi extremal proporciona un empaquetamiento físico
\(P\) de triángulos y copias de \(K_4\) en \(G\) y cantidades no
negativas \(m\) y \(A\) tales que

\[
e(G)-g(P)\le B(n,p)-\frac m{20}-\frac A2,
\qquad
B(n,p)=p(n-p)-\binom p2,
\tag{3.3}
\]

donde \(p\) es un entero entre cero y \(n\). Aquí \(m\) y \(A\) son los
recuentos efectivos de aristas exteriores y enlaces ausentes en las
cuentas físicas conservadas, respecto de la raíz regularizada. No son
errores indeterminados introducidos después de la construcción.

La calibración formal concreta utiliza \(\eta=10^{-16}\) y
\(n\ge4\cdot10^{12}\) para el Resultado de entrada 3.2. No optimizamos
estas constantes. El Resultado de entrada 3.1 puede requerir un umbral
mucho mayor.

El Resultado de entrada 3.2 combina {[}2, Proposición 4.3, Teorema 5.0 y
Lema 5.3, ecuación (5.15){]}. Su demostración tiene dos salidas
distintas. El descenso mediante copias de vértices localiza el grafo
original cerca de un grafo split completo de comparación. La
regularización de la raíz y la construcción terminal producen después el
empaquetamiento y las cuentas de (3.3), directamente en \(G\). Esta
distinción importa: el descenso se utiliza para establecer estructura,
no para transportar una sucesión de particiones de vuelta a \(G\).

\subsection{4. La deducción de Erdős
81}\label{la-deducciuxf3n-de-erdux151s-81}

Recordemos primero la cota entera para el valor de referencia split.
Completando el cuadrado se obtiene

\[
B(n,p)=\frac{(2n+1)^2}{24}
-\frac{(6p-2n-1)^2}{24}
\le\frac{n(n+1)}6+\frac1{24}.
\tag{4.1}
\]

El primer numerador es \(4n(n+1)+1\), lo que produce el término final
\(1/24\). El cuadrado sustraído es no negativo para todo entero \(p\);
la estimación no requiere, por tanto, elegir un núcleo extremizante.

El número \(B(n,p)\) es entero. La parte fraccionaria de \(n(n+1)/6\) es
múltiplo de \(1/6\), de modo que sumar \(1/24\) no puede alcanzar el
entero siguiente. En consecuencia,

\[
B(n,p)\le M(n).
\tag{4.2}
\]

\Needspace{6\baselineskip}
\textbf{Teorema 4.1.} Suponiendo los Resultados de entrada 3.1 y 3.2,
existe \(N\) tal que todo grafo cordal \(G\) de orden \(n\ge N\) admite
una partición en cliques con todas sus piezas de orden a lo sumo cuatro
y con a lo sumo \(M(n)\) piezas.

\textbf{Demostración.} Apliquemos el Resultado de entrada 3.1 con
\(\zeta=\eta/2\) y tomemos \(N\) suficientemente grande para ambos
resultados de entrada. Escribamos \(W^*=W^*(G)\). Hay dos casos. La
tolerancia se elige antes del grafo: un único umbral, el máximo de los
dos umbrales de entrada, sirve para todos los grafos considerados a
continuación.

Si \(e(G)-W^*<n^2/6-\eta n^2\), el Resultado de entrada 3.1 proporciona
un empaquetamiento físico \(P\) que satisface

\[
e(G)-g(P)
\le e(G)-W^*+\frac\eta2n^2
<\frac{n^2}{6}-\frac\eta2n^2
\le\frac{n(n+1)}6.
\tag{4.3}
\]

La primera desigualdad suma la pérdida de redondeo al coste fraccional
de completación. La desigualdad estricta del régimen lejano proporciona
el doble de la pérdida de redondeo permitida, por lo que se conserva la
mitad de ese margen. La comparación final solo utiliza \(n\ge0\).

Como \(e(G)-g(P)\) es entero, es a lo sumo \(M(n)\). Por tanto, el
certificado de una clase de la Sección 2, con \(R=e(G)\), satisface
(2.4) para \(T=M(n)\). La Proposición 2.1 da la partición requerida.

En caso contrario se cumple (3.2). El Resultado de entrada 3.2
proporciona un empaquetamiento físico con

\[
e(G)-g(P)\le B(n,p)-\frac m{20}-\frac A2
\le B(n,p)\le M(n).
\tag{4.4}
\]

La desigualdad intermedia descarta dos créditos no negativos, \(m/20\) y
\(A/2\); la última utiliza la cota entera de referencia (4.2). Las tres
estimaciones se refieren al mismo empaquetamiento proporcionado por el
Resultado de entrada 3.2. El mismo certificado y la Proposición 2.1
concluyen este caso. Todos los cliques seleccionados tienen orden tres o
cuatro, y la completación solo añade piezas \(K_2\).

\textbf{Corolario 4.2.} Existe un entero absoluto \(b\ge0\) tal que todo
grafo cordal \(G\) de cualquier orden \(n\) satisface

\[
\operatorname{cp}(G)\le c_4(G)\le M(n)+b
\le\frac{n^2}{6}+\frac n6+b.
\tag{4.5}
\]

\textbf{Demostración.} Sea \(N\) como en el Teorema 4.1 y tomemos
\(b=N^2\). Para \(n\ge N\) utilizamos el teorema. Para \(n<N\), las
aristas individuales forman una partición con a lo sumo
\(\binom n2\le N^2\) piezas. El grafo vacío también satisface la cota.
Para \(n\ge1\), la última expresión de (4.5) es a lo sumo
\(n^2/6+(b+1/6)n\); en \(n=0\) no se necesitan piezas. Esta es la cota
solicitada en el Problema 81 de Erdős.

La conclusión ya está establecida en el Paper IV; esta nota proporciona
un ensamblaje explícito de transferencia a partir de sus resultados
estructurales y de redondeo. No demuestra que \(b=0\) para todos los
órdenes ni un umbral mejorado.

\subsection{5. Qué reemplaza la
transferencia}\label{quuxe9-reemplaza-la-transferencia}

Tres preguntas aparecen en distintas etapas del argumento. La
localización pregunta por qué un grafo cordal casi extremal tiene una
raíz útil. La realización pregunta cómo convertir demandas auxiliares en
cliques del grafo original. La completación pregunta si la ganancia de
esos cliques paga el número de piezas de partición deseado.

El teorema de transferencia aborda las dos últimas preguntas una vez
proporcionadas las realizaciones y sus pérdidas. Puede reemplazar una
conversión basada en coloreabilidad cuando las clases admiten
realizaciones físicas de menor coste. Por sí solo no encuentra una raíz
ni demuestra (3.3). Por tanto, no reemplaza el papel estructural de la
copia de vértices ni de la desestabilización en esta deducción.

La primera entrada utilizada en {[}7, Sección 8{]}, y comparada con IV
en {[}2, Sección 8{]}, es un argumento de localización: se sigue un
camino de copias y se identifica su primera visita a un entorno
prescrito de la familia split. En cambio, la implementación utilizada
para el Resultado de entrada 3.2 invoca el teorema existente de
localización por descenso. Esta elección evita un selector de primera
entrada en la cadena de demostración seleccionada. Se trata de una
propiedad del resultado estructural de entrada, no de una consecuencia
de la desigualdad abstracta de transferencia. Algunos módulos fuente
también contienen el desarrollo alternativo de primera entrada;
importarlos no significa que su selector sea una premisa del teorema
utilizado aquí.

Conservar las cuentas físicas ofrece una ventaja relacionada. Si cada
paso de copia exigiera transportar una partición hacia atrás, las
pérdidas podrían acumularse a lo largo del camino. La ruta anterior
primero localiza \(G\) y después construye su empaquetamiento una sola
vez. La pérdida que paga el certificado de transferencia es la de esta
realización final. La longitud del camino de copias no aparece como un
coste adicional de partición.

\subsection{6. Lugar en la serie}\label{lugar-en-la-serie}

Los Papers I y II desarrollan el marco de perfiles y estructura
fraccional; el Paper III proporciona las herramientas de empaquetamiento
estructurado y redondeo. El Paper IV establece la localización y las
construcciones físicas requeridas aquí. El Paper V separa el enunciado
general de transferencia y su coste de realizabilidad. Esta nota hace
explícita la interfaz entre los dos últimos papers, conservando la
atribución del trabajo anterior.

La nota documenta una organización alternativa de parte de la
demostración de IV para la entrega que acompaña al Paper V. No es una
segunda demostración matemáticamente independiente de la cota extremal.
El cálculo del presupuesto y el proyecto formal separado documentan
precisamente cómo su teorema de certificados recupera la cota
establecida a partir de las construcciones de IV.

El ejemplo también delimita el alcance de una afirmación de aplicación.
Volver a expresar un buen empaquetamiento como un certificado de una
clase establece su compatibilidad con el mecanismo, pero no demuestra
que este haya aportado nuevos ahorros combinatorios. Una aplicación más
fuerte construiría realizaciones de bajo coste directamente a partir de
información estructural más débil. Esa sigue siendo una dirección de
investigación separada.

\subsection{Formalización y
reproducibilidad}\label{formalizaciuxf3n-y-reproducibilidad}

El proyecto separado en Lean 4.35 formaliza el adaptador de
transferencia física y el ensamblaje a partir de los insumos demostrados
de IV y V. Su auditoría de dependencias seleccionadas comprueba el uso
efectivo de la transferencia por clases realizables de V y de la
localización por descenso de IV, y excluye el teorema final de cota
global de IV. Esta exclusión describe el ensamblaje, no una
independencia matemática respecto de IV. El companion formal identifica
el paquete final congelado de fuentes, la auditoría completa de axiomas
locales, las interfaces públicas y los comandos de reproducción.
\href{https://github.com/jtraverso/erdos-81-chordal-clique-partitions/tree/main/series_notes/NOTE_I_ERDOS81_TRANSFER/05_formalization}{Fuentes
y companion formal}.

\subsection{Agradecimientos}\label{agradecimientos}

El autor agradece profundamente a su esposa María Paz y a sus hijos
Lucas, Juan Cristóbal, Francisca, Raimundo y Benjamín su amor, paciencia
y apoyo.

\subsection{Uso de inteligencia artificial y herramientas
computacionales}\label{uso-de-inteligencia-artificial-y-herramientas-computacionales}

En los artículos fuente utilizados aquí, se utilizaron Claude, de
Anthropic, y ChatGPT/Codex, de OpenAI, para explorar argumentos,
comprobarlos y preparar los manuscritos. Aristotle, de Harmonic,
participó en la exploración de caminos, la búsqueda de contraejemplos,
la elaboración de pruebas y su formalización y revisión en Lean; no sólo
tradujo pruebas terminadas.

Para esta nota, ChatGPT/Codex ayudó a ensamblar el argumento de las
fuentes, construir y comprobar los adaptadores Lean y preparar la
exposición y la auditoría interna.

Certo {[}5{]} y Jacobian {[}6{]} son las herramientas computacionales
usadas en el programa de investigación subyacente para explorar
instancias, comprobar identidades y evaluar candidatos, según se declara
en el Paper IV. Su uso se declara por separado de las dependencias
lógicas de la prueba formal. El companion separado y los registros de
auditoría identifican las comprobaciones realizadas específicamente para
esta nota, incluidas sus pruebas de regresión en Python. Los cálculos se
revisaron según su alcance; una búsqueda finita no reemplaza una prueba
universal.

\Needspace{3\baselineskip}El autor conserva la responsabilidad por los argumentos, las citas, el
código y la presentación. Ningún sistema de IA figura como autor.

\subsection{Referencias}\small\label{referencias}

{[}1{]} P. Erdős, E. T. Ordman y Y. Zalcstein, ``Clique partitions of
chordal graphs,'' Combinatorics, Probability and Computing 2 (1993),
409--415. \url{https://doi.org/10.1017/S0963548300000808}.

{[}2{]} J. P. Traverso Gianini, Clique partitions with rooted simplicial
defect: quantitative stability and sharp eventual bounds, Paper IV,
preprint v1.23, 1 de octubre de 2026. La versión conservada en la
distribución de la serie es la fuente de los Resultados de entrada 3.1 y
3.2.
\href{https://github.com/jtraverso/erdos-81-chordal-clique-partitions/blob/054593443bdece587dda9686a411bb18b0ca34c5/preprints/PAPER_IV/01_manuscript/PAPER_IV_preprint_v1.23_en.md}{Manuscrito
y materiales}. Depósito:
\href{https://doi.org/10.5281/zenodo.23089131}{doi:10.5281/zenodo.23089131}.

{[}3{]} J. P. Traverso Gianini, Colorability and Realizability Defects
in Resource-Dominated Packing Relaxations, Paper V, \mbox{Preprint 1.0} y
paquete formal de fuentes complementario, 18 de octubre de 2026. El
teorema de transferencia por clases realizables es el resultado
utilizado en la Sección 2.
\href{https://github.com/jtraverso/erdos-81-chordal-clique-partitions/tree/main/preprints/PAPER_V}{Manuscrito
y materiales}.

{[}4{]} T. F. Bloom, ``Erdős Problem \#81,'' Erdős Problems,
https://www.erdosproblems.com/81. Se cita para identificar el problema;
esta nota no certifica el estado actual de revisión externa de las
soluciones propuestas.

{[}5{]} J. P. Traverso Gianini, \emph{Certo: herramientas
computacionales con certificados verificables}.
\href{https://github.com/jtraverso/certo-math}{Manuscrito y materiales}.

{[}6{]} Morluto y colaboradores, \emph{Jacobian: herramientas
matemáticas componibles}.
\href{https://github.com/morluto/jacobian}{Manuscrito y materiales}.

{[}7{]} Anonymous, \emph{Clique partitions of chordal graphs with linear
error}, preprint, 8 de septiembre de 2026. Versión comparada en el Paper
IV: commit \nolinkurl{cbde8a0a0563372b23b1b39a44180d2c0fb02f44}.
\href{https://github.com/N0zoM1z0/erdos-81/blob/cbde8a0a0563372b23b1b39a44180d2c0fb02f44/manuscript/main.pdf}{Manuscrito
y materiales}. Se conserva la firma del manuscrito. El README de ese
commit atribuye la demostración matemática a Morluto y el manuscrito y
la implementación Lean a N0zoM1z0, con Jacobian de Preference Labs como
parte del proyecto; estos créditos se distinguen de la autoría impresa.

{[}8{]} O. Okechukwu, \emph{Clique partitions and bounded simplicial
defect}, arXiv:2609.20871v1, 15 de septiembre de 2026.
\url{https://arxiv.org/abs/2609.20871v1}. La comparación y el alcance
utilizados aquí son los documentados en el Paper IV, Sección 8.1.

{[}9{]} P. E. Haxell y V. Rödl, ``Integer and fractional packings in
dense graphs,'' \emph{Combinatorica} \textbf{21} (2001), 13--38.
\url{https://doi.org/10.1007/s004930170003}.

{[}10{]} R. Yuster, ``Integer and fractional packing of families of
graphs,'' \emph{Random Structures \& Algorithms} \textbf{26} (2005),
110--118. \url{https://doi.org/10.1002/rsa.20048}.

{[}11{]} D. Rohatgi, J. C. Urschel, y J. Wellens, ``Regarding Two
Conjectures on Clique and Biclique Partitions,'' \emph{Electronic
Journal of Combinatorics} \textbf{28}(4) (2021).
\url{https://doi.org/10.37236/9564}. Versión de los autores:
\url{https://arxiv.org/abs/2005.02529}.

\end{document}
